Synthetic generators of equity-return time series produce plausible sample paths that match the statistical structure of real markets without replaying the single realized history, supporting stress testing, regulatory backtesting, and scenario design against conditions absent from the observed record~\cite{assefa2020generating}. A useful generator should reproduce the three symmetric stylized facts of daily returns at once: heavy-tailed marginals, negligible linear autocorrelation, and slow decay of the autocorrelation function (ACF) of absolute returns~\cite{mandelbrot1963variation, cont2001empirical}. Generalised autoregressive conditional heteroskedasticity (GARCH) models~\cite{engle1982autoregressive, bollerslev1986generalized} capture volatility clustering but not the multi-regime structure of the marginal; i.i.d.\ generators~\cite{efron1979bootstrap} match the marginal but lose temporal persistence; deep generators~\cite{yoon2019timegan, wiese2020quantgan} reproduce some stylized facts at the cost of interpretability and reproducibility. One negative result has shaped two decades of work on regime-switching models. Ryd\'{e}n et al.~\cite{ryden1998stylized} fitted Gaussian-emission hidden Markov models (HMMs) at low state counts and concluded that the slow decay of the absolute-return ACF was the one stylized fact the model could not generate, ``at least not easily by ML estimation'': a finite-state chain produces only exponential ACF decay, and the joint fit, forced to match the remaining properties of the data, did not land on a slow one. Bulla and Bulla~\cite{bulla2006stylized} attributed the failure to the geometric sojourn-time distribution of the standard Markov chain and recovered the ACF fit with a hidden semi-Markov chain, at the cost of much greater complexity. We separate the low-$K$ Gaussian limitation into two possible failure channels: a temporal one (the absolute growth-rate ACF must allow slow-decay modes) and a distributional one (the marginal mixture must carry enough components to match observed excess kurtosis). Ryd\'{e}n's own outlier-reduced subseries broke on the temporal channel; on our non-outlier-reduced data the distributional channel binds instead. A standard closed-form identity for the absolute growth-rate ACF, written as a sum over the non-unit eigenvalues of the transition matrix~\cite{hamilton1994time, timmermann2000moments}, bounds the number of decay modes by the rank of the centred transition matrix. We ask whether that bound actually matters on equity-return data once enough latent states are allowed.

We answer this with an established model class, the continuous-emission hidden Markov model (CHMM): a hidden chain of market regimes in which each regime emits the day's growth rate from its own probability distribution. The regime chain controls the geometric modes available to the absolute growth-rate ACF, while the per-regime distributions control the heavy-tailed marginal, allowing the temporal and distributional constraints of the low-state Gaussian fit to be examined separately. Our contribution is not the CHMM itself. It is a unified comparison of Gaussian, Student-$t$, Laplace, and generalized-error emissions under the same forward--backward framework; quantile-based initialization, which seeds each regime with its own slice of the observed distribution so the fit converges with every regime still in use; an empirical spectral diagnosis of when the finite-mode bound is active; and a regime-conditional Value-at-Risk (VaR) head that forecasts tail risk from the model's running estimate of which regime the day belongs to. Against deep generative baselines, the fitted model competes at a parameter budget small enough to inspect: 12 free parameters at three states, against the orders-of-magnitude larger convolutional generator--critic architecture of QuantGAN~\cite{wiese2020quantgan}, with each per-state location, scale, and shape directly readable as an economic regime.

We evaluate on a six-fold rolling-origin walk-forward over a decade of SPY and on a sector-balanced 30-ticker panel across ten Global Industry Classification Standard (GICS) sectors, reducing reliance on a single window or asset universe. With enough latent states and a heavy-tailed emission, the model reproduces all three symmetric Cont stylized facts and narrows the excess-kurtosis gap left by the Gaussian baseline. Under the diagnostics used here, the rank bound from the ACF identity is not empirically active at the typical ticker, although a finite-state HMM still permits only finitely many geometric decay modes and the result does not establish a general power-law or multi-scale approximation. The regime-conditional VaR is not rejected by the Christoffersen joint conditional-coverage test~\cite{christoffersen1998evaluating} on the main window or on most walk-forward folds. We therefore interpret the Ryd\'{e}n low-state-count failure more narrowly: on these daily US-equity data and metrics, increasing marginal flexibility resolves more of the observed fit gap than adding ACF modes beyond the selected state count. An i.i.d.\ bootstrap and a maximum-likelihood hidden semi-Markov model (HSMM) beat the CHMM on raw single-window fit, but the fitted CHMM additionally supports the regime-conditional risk forecast exercised here. We make no formal privacy guarantee for its synthetic series. The scope is daily US equities: a non-equity stress test on the GLD exchange-traded fund (ETF) broke down under static fitting, and we recommend periodic refit within that scope.
